\section{Kết luận}
Trải qua quá trình nghiên cứu, phân tích và triển khai thực tế, đề tài ``Xây dựng website thương mại điện tử thời trang tích hợp Virtual Try-On AI và thử đồ thực tế tăng cường'' đã hoàn thành đúng tiến độ và đạt được các mục tiêu trọng tâm đề ra tại Chương 1. Đề tài đã xây dựng thành công một nền tảng mua sắm trực tuyến kết hợp trí tuệ nhân tạo và công nghệ thực tế tăng cường, mở ra hướng tiếp cận mới cho trải nghiệm thời trang số.

Cụ thể, nhóm đã đạt được các thành tựu sau:
\begin{itemize}[noitemsep]
    \item \textbf{Về mặt hạ tầng và công nghệ:} Xây dựng thành công hệ thống Full-stack hoàn chỉnh với Nuxt.js 3 (SPA mode) cho Frontend, FastAPI cho Backend và kiến trúc cơ sở dữ liệu kép gồm Firebase Firestore (lưu trữ sản phẩm, đơn hàng) và Neon PostgreSQL với pgvector (dữ liệu người dùng, AI, vector search). Hệ thống được tích hợp Cloudinary để quản lý tài nguyên media và Stripe làm cổng thanh toán quốc tế.

    \item \textbf{Về nghiệp vụ thương mại điện tử:} Hệ thống đã triển khai hoàn chỉnh quy trình mua sắm trực tuyến hiện đại, bao gồm duyệt sản phẩm theo danh mục, quản lý giỏ hàng, hệ thống voucher giảm giá, và tích hợp thanh toán qua Stripe. Người dùng có thể xác thực bằng Google Sign-In thông qua Firebase Auth, quản lý nhiều địa chỉ giao hàng và theo dõi trạng thái đơn hàng theo thời gian thực.

    \item \textbf{Về đột phá công nghệ AI Virtual Try-On:} Đây là điểm cốt lõi của đề tài. Nhóm đã tích hợp thành công mô hình FitDiT~\cite{fitdit_paper} --- một Diffusion Transformer tiên tiến dựa trên nền tảng Stable Diffusion 3 --- cho phép người dùng upload ảnh cá nhân và xem trước trang phục ảo với độ phân giải lên đến 1536$\times$2048. Pipeline xử lý tự động bao gồm DWpose để phát hiện keypoint cơ thể và Human Parsing để tách vùng trang phục, đảm bảo kết quả thử đồ chân thực và tự nhiên.

    \item \textbf{Về thử đồ thực tế tăng cường (AR):} Ứng dụng tích hợp Snap Camera Kit~\cite{snap_camera_kit_docs} kết hợp Lens Studio~\cite{lens_studio_docs} với công nghệ ObjectTracking3D~\cite{snap_body_tracking_docs} cho phép người dùng thử đồ trực tiếp qua camera theo thời gian thực ngay trên trình duyệt web. Hệ thống carousel 3D cho phép chuyển đổi linh hoạt giữa các mẫu trang phục thông qua thao tác cử chỉ tay.

    \item \textbf{Về gợi ý size thông minh:} Hệ thống Fit Assessment dựa trên phương pháp Ease Allowance tự động so sánh số đo cơ thể người dùng với thông số kỹ thuật từng size của trang phục, đưa ra khuyến nghị size phù hợp kèm cảnh báo các vùng có nguy cơ không vừa (ngực, eo, hông).

    \item \textbf{Về chatbot tư vấn thời trang:} Chatbot AI sử dụng kiến trúc RAG (Retrieval-Augmented Generation)~\cite{rag_lewis_2020} kết hợp mô hình ngôn ngữ Ollama qwen2.5:7b~\cite{ollama_docs} chạy cục bộ và pgvector~\cite{pgvector_github} để tìm kiếm ngữ nghĩa, đóng vai trò tư vấn viên thời trang ảo, hỗ trợ người dùng lựa chọn sản phẩm phù hợp.

    \item \textbf{Về giá trị thực tiễn:} Hệ thống đã giải quyết bài toán thực tiễn trong thương mại điện tử thời trang --- giúp khách hàng giảm thiểu rủi ro mua sai size hoặc không ưng sản phẩm khi mua online, đồng thời cung cấp công cụ quản lý cửa hàng toàn diện cho người bán với dashboard theo dõi đơn hàng, sản phẩm và voucher tập trung.
\end{itemize}